\section{Physical-validation details}
\subsection{Board harness and feasibility amendments}

The board study uses paired SFT/RF candidates from the same design and applies
the same multiplicity-aware selection rule to both policies.
Fig.~\ref{fig:board-architecture} shows this selection contract, the implemented
test harness, and the trace-validation path. All selected DUTs, an echo canary,
and the selector interface occupy one bitstream/HWH pair. Each complete run
sweeps the same clock grid, recomputes the passing boundary from raw scores,
and records the bitstream, HWH, selector-map, and selection-manifest hashes.
This common harness allows the paired outputs to be measured under the same
clock and capture conditions.

\begin{figure*}[!t]
\centering
\includegraphics[width=\textwidth]{fig04_board_architecture_v2}
\caption{Descriptive board-validation architecture.
(a) The amended correctness-eligible selection uses the same multiplicity and
hash rules for SFT and RF, with no frequency-based reselection.
(b) The processing system configures the clock and accesses the capture
buffer through memory-mapped AXI. Counter-driven SFT/RF DUTs and the echo
canary coexist in one programmable-logic image; the selector chooses the
captured output. Solid paths carry data or execution flow; dashed paths show
clock and capture control. Logical buses are simplified for readability.
(c) Captured traces are compared with per-DUT golden streams, and clock
boundaries are accepted only with the specified canary headroom.}
\label{fig:board-architecture}
\end{figure*}

The first outcome-independent family hash selected a reverse-FIR design with no
correct candidate in either paired arm, so that attempt stopped without a
silicon outcome.  Before any board measurement, eligibility was amended to
require a correct candidate in both arms and the same outcome-independent hash
was reapplied.  The common sweep floor was also lowered after sealed Vivado
rows showed that the original range would not bracket every selected DUT.  No
reward, timing, area, power, or resource outcome was used to choose the
replacement.  Nevertheless, these feasibility amendments make the resulting
board subset descriptive rather than independently confirmatory.

\subsection{Selected setup-closure pilot}
% AUTO-GENERATED by paper/make_revision_evidence.py.
\begin{table}[t]
\centering
\caption{Descriptive setup-closure sensitivity (MHz). Original: frozen WNS-derived value. Setup: conservative passing endpoint under the pilot flow.}
\label{tab:setup-closure}
\small
\begin{tabular}{@{}lrrrr@{}}
\toprule
Design & SFT orig. & SFT setup & RF orig. & RF setup \\
\midrule
fir42\_v6\_8b & \revisionclaim{ClosureASftOriginal} & \revisionclaim{ClosureASftSetup} & \revisionclaim{ClosureARfOriginal} & \revisionclaim{ClosureARfSetup} \\
firr20\_v6 & \revisionclaim{ClosureBSftOriginal} & \revisionclaim{ClosureBSftSetup} & \revisionclaim{ClosureBRfOriginal} & \revisionclaim{ClosureBRfSetup} \\
iir18\_v7 & \revisionclaim{ClosureCSftOriginal} & \revisionclaim{ClosureCSftSetup} & \revisionclaim{ClosureCRfOriginal} & \revisionclaim{ClosureCRfSetup} \\
med21 & \revisionclaim{ClosureDSftOriginal} & \revisionclaim{ClosureDSftSetup} & \revisionclaim{ClosureDRfOriginal} & \revisionclaim{ClosureDRfSetup} \\
poly6\_v13\_8b & \revisionclaim{ClosureESftOriginal} & \revisionclaim{ClosureESftSetup} & \revisionclaim{ClosureERfOriginal} & \revisionclaim{ClosureERfSetup} \\
\bottomrule
\end{tabular}

\end{table}


A separate pilot evaluates \revisionclaim{ClosureCandidateCount} selected
candidates under Vivado \revisionclaim{ClosureVersion}, applying constraints
before synthesis and bracketing a setup-passing period for each candidate.
It therefore changes both the tool version and constraint handling relative
to the primary flow. The selected candidates yield Spearman correlation
\revisionclaim{ClosureRho} between original timing-derived and setup-passing
frequencies, with median symmetric percentage error
\revisionclaim{ClosureMedianSmape}\%.  None of the
\revisionclaim{ClosureNonTiedPairs} non-tied pairs reverses direction
(Table~\ref{tab:setup-closure}).  However, the largest candidate error is
\revisionclaim{ClosureWorstSmape}\%: the reverse-FIR RF candidate changes from
\revisionclaim{ClosureBRfOriginal} to
\revisionclaim{ClosureBRfSetup}~MHz.  The median error therefore must not be
read as a uniform accuracy guarantee.  The check supports relative ordering
on this subset, while absolute frequency remains sensitive to implementation
conditions.  Setup passing alone is not full timing sign-off.

\subsection{Complete cross-version coverage diagnostic}
% AUTO-GENERATED by paper/make_revision_evidence.py.
\begin{table}[t]
\centering
\caption{Cross-version coverage and timing agreement. Error is median symmetric percentage error on common valid candidates. Unresolved cases remain inconclusive.}
\label{tab:coverage}
\small
\begin{tabular}{@{}lrrrrr@{}}
\toprule
Arm & Total & Valid & No path & Inconcl. & Error (\%) \\
\midrule
SFT & \revisionclaim{CoverageSftCount} & \revisionclaim{CoverageSftValid} & \revisionclaim{CoverageSftInvalid} & \revisionclaim{CoverageSftInconclusive} & \revisionclaim{CoverageSftMedianError} \\
RF seed one & \revisionclaim{CoverageRfOneCount} & \revisionclaim{CoverageRfOneValid} & \revisionclaim{CoverageRfOneInvalid} & \revisionclaim{CoverageRfOneInconclusive} & \revisionclaim{CoverageRfOneMedianError} \\
RF seed two & \revisionclaim{CoverageRfTwoCount} & \revisionclaim{CoverageRfTwoValid} & \revisionclaim{CoverageRfTwoInvalid} & \revisionclaim{CoverageRfTwoInconclusive} & \revisionclaim{CoverageRfTwoMedianError} \\
\bottomrule
\end{tabular}

\end{table}

The archived diagnostic preserves RTL, device, clock request, and constraint
order while changing the Vivado version. It covers
\revisionclaim{CoverageCandidates} original SFT/RF candidates. All
\revisionclaim{CoverageValid} setup-path-present candidates have zero counts in
the six named constraint-coverage checks. The confirmed missing-path case is
\revisionclaim{CoverageInvalidModule}; \revisionclaim{CoverageInconclusive}
other implementations remain inconclusive after the retained attempts.
\revisionclaim{CoverageWarningCount} SFT candidates also have multi-driver
synthesis warnings; these are recorded without silently converting them into
physical failures.

Median common-valid symmetric errors appear in Table~\ref{tab:coverage}; the
maximum is \revisionclaim{CoverageMaxError}\%. The complete RF seed-two
diagnostic mean is \revisionclaim{CoverageRfTwoReplay}~MHz versus
\revisionclaim{CoverageRfTwoOriginal}~MHz originally. SFT and RF seed-one means
are unavailable because of the inconclusive cases, which are not assigned zero.
Missing-path replays receive zero and all original draw denominators remain.
These results describe timing sensitivity and coverage, not period-bracketed
closure, synthesized functional equivalence, full timing sign-off, or a corrected
primary confidence interval. The retained primary values are unchanged.
